\documentclass[12pt,a4paper]{article}

% =========================
% Packages
% =========================
\usepackage[margin=0.85in]{geometry}
\usepackage[T1]{fontenc}
\usepackage[utf8]{inputenc}
\usepackage{lmodern}
\usepackage{microtype}
\usepackage{hyperref}
\usepackage{xcolor}
\usepackage{enumitem}
\usepackage{booktabs}
\usepackage{longtable}
\usepackage{tabularx}
\usepackage{array}
\usepackage{tcolorbox}
\usepackage{titlesec}
\usepackage{listings}
\usepackage{pifont}

% =========================
% Styling
% =========================
\definecolor{primary}{RGB}{255,102,0} % Swedbank Orange
\definecolor{secondary}{RGB}{45,45,45}
\definecolor{codebg}{RGB}{245,245,245}
\definecolor{lightblue}{RGB}{235,245,255}
\definecolor{lightred}{RGB}{255,240,240}
\definecolor{lightgreen}{RGB}{240,255,240}
\definecolor{lightyellow}{RGB}{255,250,230}

\hypersetup{
    colorlinks=true,
    linkcolor=primary,
    urlcolor=primary,
    pdftitle={Detailed 7-Day AML Graph-Fraud Pipeline Plan},
    pdfauthor={Portfolio Project Plan}
}

\titleformat{\section}
{\Large\bfseries\color{primary}}
{\thesection}{1em}{}

\titleformat{\subsection}
{\large\bfseries\color{secondary}}
{\thesubsection}{1em}{}

\titleformat{\subsubsection}
{\normalsize\bfseries}
{\thesubsubsection}{1em}{}

\setlist[itemize]{leftmargin=*, itemsep=0.2em}
\setlist[enumerate]{leftmargin=*, itemsep=0.2em}

\lstset{
    backgroundcolor=\color{codebg},
    basicstyle=\ttfamily\small,
    breaklines=true,
    frame=single,
    rulecolor=\color{black!20},
    columns=fullflexible,
    keepspaces=true
}

\newcommand{\done}{\ding{51}}
\newcommand{\todo}{\ding{113}}

% =========================
% Title
% =========================
\title{
    \textbf{\Huge Detailed 7-Day Execution Blueprint}\\
    \vspace{0.4em}
    \Large Real-Time AML Graph-Fraud Pipeline\\
    \vspace{0.3em}
    \large Master's Level Portfolio Project for Swedbank
}
\author{}
\date{August 2026}

\begin{document}

\maketitle

\begin{abstract}
\noindent
This document provides a day-by-day, granular execution plan to build, train, and deploy an event-driven Graph Neural Network (GraphSAGE) for Anti-Money Laundering (AML) within a strict 7-day timeline. It details specific algorithmic tasks, file structures, and MLOps integrations required to prove senior-level competency in financial machine learning.
\end{abstract}

\vspace{0.5em}
\hrule
\vspace{1em}

\tableofcontents
\newpage

% ============================================================
\section{Complete Repository Structure}

The following directory tree encompasses all files generated across the 7-day sprint. It keeps data ingestion, model definition, training loops, and streaming services strictly modular—a structure expected in enterprise MLOps.

\begin{lstlisting}[language=bash]
aml-graph-pipeline/
|-- README.md                   # Final documentation and run instructions (Day 7)
|-- docker-compose.yml          # Kafka & MLflow services config (Day 5)
|-- requirements.txt            # Python dependencies
|
|-- data/
|   `-- raw/                    # Downloaded Elliptic Bitcoin Dataset
|
|-- mlruns/                     # Local SQLite DB for MLflow tracking (Auto-generated Day 4)
|
|-- src/
|   |-- data/
|   |   |-- loader.py           # PyG dataset ingestion & temporal splits (Day 1)
|   |   `-- baseline_xgboost.py # Tabular-only baseline model (Day 1)
|   |
|   |-- model/
|   |   |-- graphsage.py        # 2-layer GraphSAGE PyTorch architecture (Day 2)
|   |   `-- ensemble.py         # Custom MLflow PyFunc wrapping PyG + XGBoost (Day 4)
|   |
|   |-- training/
|   |   |-- train_gnn.py        # Pre-training script for GraphSAGE (Day 2)
|   |   |-- hybrid_train.py     # Training XGBoost on embeddings + raw features (Day 3)
|   |   `-- train_pipeline.py   # Final MLflow training loop & registry logging (Day 4)
|   |
|   `-- streaming/
|       |-- producer.py         # Simulates live Kafka transaction events (Day 5)
|       `-- consumer.py         # Real-time inductive inference & SHAP alerts (Day 6)
|
`-- docs/
    `-- architecture.mermaid    # System design diagram for Swedbank pitch (Day 7)
\end{lstlisting}

% ============================================================
\section{The 7-Day Sprint Plan}

This schedule is aggressive but achievable. It separates the deep learning research (Days 1--3) from the MLOps engineering (Days 4--6), culminating in presentation prep (Day 7).

% ------------------------------------------------------------
\subsection{Day 1: Graph Ingestion \& Baseline Establishment}
\textbf{Objective:} Load the dataset, enforce temporal constraints, and prove why a GNN is necessary by setting a tabular baseline.

\begin{itemize}
    \item[\todo] \textbf{Scripting:} Create \texttt{src/data/loader.py}. Use \texttt{torch\_geometric.datasets.EllipticBitcoinDataset} to download the graph.
    \item[\todo] \textbf{Temporal Split:} Implement strict chronological splitting. Train: timesteps 1--30. Validation: 31--34. Test: 35--49.
    \item[\todo] \textbf{Baseline Model:} In \texttt{src/data/baseline\_xgboost.py}, extract the raw node features (166 tabular columns) for the training set. Train an out-of-the-box XGBoost classifier \textit{without} graph features. 
    \item[\todo] \textbf{Metrics:} Calculate the baseline PR-AUC (Precision-Recall Area Under Curve). This is the score your GNN must beat.
\end{itemize}

% ------------------------------------------------------------
\subsection{Day 2: Inductive GraphSAGE Architecture}
\textbf{Objective:} Build the neural network that understands transaction topologies.

\begin{itemize}
    \item[\todo] \textbf{Sampling Setup:} Implement PyG's \texttt{NeighborLoader} in \texttt{src/training/train\_gnn.py}. Configure it to sample 25 neighbors at hop 1, and 10 neighbors at hop 2.
    \item[\todo] \textbf{Architecture:} Create \texttt{src/model/graphsage.py}. Subclass \texttt{torch.nn.Module}. Create two \texttt{SAGEConv} layers.
    \item[\todo] \textbf{Forward Pass:} Design the \texttt{forward()} method to output a dense 64-dimensional embedding, \textit{not} a softmax classification. 
    \item[\todo] \textbf{Pre-training:} Train the GraphSAGE model on the Train split using a standard Cross-Entropy loss to predict illicit nodes. Save the model weights locally.
\end{itemize}

% ------------------------------------------------------------
\subsection{Day 3: The Hybrid Ensemble (GraphSAGE + XGBoost)}
\textbf{Objective:} Fuse spatial network features with raw tabular features for final classification.

\begin{itemize}
    \item[\todo] \textbf{Feature Extraction:} In \texttt{src/training/hybrid\_train.py}, pass all Training nodes through the frozen GraphSAGE model to extract their 64d embeddings.
    \item[\todo] \textbf{Concatenation:} Merge the 166 original features with the 64 graph embeddings (resulting in 230 features).
    \item[\todo] \textbf{Ensemble Training:} Train a new XGBoost classifier on this 230-feature dataset.
    \item[\todo] \textbf{Validation:} Compare this hybrid model's PR-AUC against the Day 1 baseline. You should see a noticeable jump in recall for illicit transactions.
\end{itemize}

% ------------------------------------------------------------
\subsection{Day 4: MLflow Tracking \& Packaging}
\textbf{Objective:} Professionalize the training loop using enterprise experiment tracking.

\begin{itemize}
    \item[\todo] \textbf{MLflow Setup:} Spin up a local MLflow tracking server using SQLite.
    \item[\todo] \textbf{Tracking Script:} Combine Days 2 and 3 into \texttt{src/training/train\_pipeline.py}. Include \texttt{mlflow.start\_run()} to log hyperparameters (learning rate, tree depth, neighbor sample size).
    \item[\todo] \textbf{Custom PyFunc:} Create \texttt{src/model/ensemble.py} holding a custom MLflow Python model class that wraps \textit{both} the PyTorch GraphSAGE weights and the XGBoost booster. 
    \item[\todo] \textbf{Artifact Logging:} Save the combined model to the MLflow model registry so the Kafka consumer only needs to load one unified artifact.
\end{itemize}

% ------------------------------------------------------------
\subsection{Day 5: Event-Driven Infrastructure (Kafka)}
\textbf{Objective:} Simulate the Swedbank real-time transaction environment.

\begin{itemize}
    \item[\todo] \textbf{Dockerization:} Write \texttt{docker-compose.yml} to launch a single-node Apache Kafka instance (using KRaft, bypassing Zookeeper for speed) alongside the MLflow server.
    \item[\todo] \textbf{Producer Script:} Create \texttt{src/streaming/producer.py}. Iterate through the Test split (timesteps 35--49) and emit each transaction as a JSON payload to a Kafka topic named \texttt{aml-transactions}.
    \item[\todo] \textbf{Throughput Testing:} Run the producer and use Kafka CLI tools to verify messages are landing in the topic at realistic speeds.
\end{itemize}

% ------------------------------------------------------------
\subsection{Day 6: Real-Time Streaming Inference}
\textbf{Objective:} Close the loop by consuming events and scoring them dynamically.

\begin{itemize}
    \item[\todo] \textbf{Consumer Script:} Create \texttt{src/streaming/consumer.py}.
    \item[\todo] \textbf{Model Loading:} Pull the latest version of the custom PyFunc model from MLflow.
    \item[\todo] \textbf{Inductive Scoring:} For every consumed Kafka message, dynamically query the local graph state for the node's neighbors, pass the subgraph through GraphSAGE, concatenate, and predict via XGBoost.
    \item[\todo] \textbf{Alerting \& Explainability:} If prediction probability $> 0.85$, print an AML Alert to the console, utilizing \texttt{shap.TreeExplainer} to log the top 3 features driving the fraud score.
\end{itemize}

% ------------------------------------------------------------
\subsection{Day 7: Documentation \& Pitch Formulation}
\textbf{Objective:} Package the code so a Swedbank engineering manager can understand it in 60 seconds.

\begin{itemize}
    \item[\todo] \textbf{README.md:} Write a comprehensive guide. Include commands to run the docker-compose and the producer/consumer.
    \item[\todo] \textbf{Architecture Diagram:} Use Mermaid.js (\texttt{docs/architecture.mermaid}) to map the flow: \textit{Data $\rightarrow$ Producer $\rightarrow$ Kafka $\rightarrow$ Consumer (PyG + XGBoost) $\rightarrow$ Alert}.
    \item[\todo] \textbf{The Pitch:} Draft a 2-minute narrative explaining \textit{why} you chose inductive GNNs over standard models for AML compliance.
\end{itemize}

\end{document}